% from ega4-2-fr.tex
\oldpage[IV-2]{114}

\begin{remarks}[5.9.12]
\label{IV.2.s5.9.12-fr}
\begin{enumerate}
\item[(i)] Soit $\mathcal C(X)$ la catégorie des $\mathcal O_X$-Modules, et
soit $\mathcal C_Z(X)$ la sous-catégorie de $\mathcal C(X)$ formée des
$\mathcal O_X$-Modules de support contenu dans $Z$; cette sous-catégorie est
localisante au sens de Gabriel, et le foncteur $\mathcal H^0_{X/Z}$ n'est
autre que le foncteur localisation de Gabriel (cf. [27]; cela fournirait une
autre démonstration de (5.9.11)). Lorsque $Z$ est fermé, le foncteur
$i^*:\mathcal C(X)\to\mathcal C(X-Z)$ (où $i:X-Z\to X$ est l'injection
canonique) définit une équivalence de catégories
\[
\mathcal C(X)/\mathcal C_Z(X)\approx\mathcal C(X-Z).
\]
\item[(ii)] Il résulte de (5.9.11) que la condition
$\mathcal H^0_{X/Z}(\mathcal F)=0$ est équivalente à
$\operatorname{Supp}(\mathcal F)\subset Z$. Elle entraîne en effet cette
dernière puisque le noyau de $\rho_{X/Z}$ est alors égal à $\mathcal F$.
Inversement, si $\operatorname{Supp}(\mathcal F)\subset Z$, il suffit
d'appliquer la seconde partie de (5.9.11) à l'unique homomorphisme
$u:\mathcal F\to0$ pour en conclure que l'homomorphisme correspondant
$v:\mathcal H^0_{X/Z}(\mathcal F)\to0$ est un isomorphisme.
\item[(iii)] Les développements précédents gardent un sens pour tout espace
annelé localement noethérien dont toute partie fermée irréductible admet
exactement un point générique. En particulier ils s'appliquent, sur un
préschéma localement noethérien, à des \emph{faisceaux de groupes abéliens}
quelconques (considérés comme Modules sur le faisceau simple associé au
préfaisceau constant $\mathbf Z$). On a encore pour tout $x\in X$
l'isomorphisme canonique (5.9.6.1), dans lequel
$\widetilde{\mathcal F}_x$ désigne le faisceau induit sur le sous-espace
$X_x$ de $X$ par le faisceau $\mathcal F$; la démonstration directe résulte
aussitôt de la définition (5.9.1.2) et du théorème de la double limite
inductive.
\end{enumerate}
\end{remarks}

\subsection{\texorpdfstring{Propriété $(S_2)$ et $Z$-clôture}{Propriété (S\string_2) et Z-clôture}}

\begin{env}[5.10.1]
\label{IV.2.s5.10.1-fr}
Soient $X$ un préschéma localement noethérien, $\mathcal F$ un
$\mathcal O_X$-Module cohérent; pour toute partie $T$ de $X$, nous poserons
\begin{equation}
\operatorname{prof}_T(\mathcal F)
=\inf_{x\in T}\operatorname{prof}(\mathcal F_x).
\tag{5.10.1.1}\label{IV.2.s5.10.1.1-fr}
\end{equation}
\end{env}

\begin{proposition}[5.10.2]
\EGAReaderTarget{ega4.c04.prop.5-10-2}{IV.2.s5.10.2-fr}
Soient $X$ un préschéma localement noethérien, $Z$ une partie de $X$ stable
par spécialisation, $\mathcal F$ un $\mathcal O_X$-Module quasi-cohérent.
Les conditions suivantes sont équivalentes:
\begin{enumerate}
\item[a)] $\mathcal F$ est $Z$-pur.
\item[b)] $\operatorname{Ass}(\mathcal F)$ ne rencontre pas $Z$.
\end{enumerate}
Si en outre $\mathcal F$ est cohérent, ces conditions sont aussi équivalentes
à la suivante:
\begin{enumerate}
\item[c)] $\operatorname{prof}_Z(\mathcal F)\geq1$.
\end{enumerate}
\end{proposition}

Dire que $\mathcal F$ est $Z$-pur signifie que pour tout ouvert noethérien
$V$ de $X$, et tout ouvert $U\supset X-Z$, l'homomorphisme de restriction
\[
\Gamma(V,\mathcal F)\longrightarrow\Gamma(V\cap U,\mathcal F)
\]
est injectif (5.9.8); mais d'après (3.1.8) cela équivaut à
$V\cap\operatorname{Ass}(\mathcal F)\subset U$, d'où l'équivalence de a) et
b). En outre, dire que $x\in\operatorname{Ass}(\mathcal F)$ signifie
qu'aucun élément de $\mathfrak m_x$ n'est $\mathcal F_x$-régulier (3.1.2),
donc, lorsque $\mathcal F$ est cohérent, cela s'écrit encore
$\operatorname{prof}(\mathcal F_x)=0$; on en déduit aussitôt dans ce cas
l'équivalence de b) et c).

\begin{corollary}[5.10.3]
\label{IV.2.s5.10.3-fr}
Soit
\[
0\longrightarrow\mathcal F'\longrightarrow\mathcal F
\longrightarrow\mathcal F''\longrightarrow0
\]
une suite exacte de $\mathcal O_X$-Modules quasi-cohérents. Si $\mathcal F$
est $Z$-pur, il en est de même de $\mathcal F'$; inversement, si
$\mathcal F'$ et $\mathcal F''$ sont $Z$-purs, il en est de même de
$\mathcal F$.
\end{corollary}

